\chapter{Timed Automata, Regions, and Zones}
\label{sec:p09-timed-automata-regions-and-zones}

Timed automata combine a finite control graph with real-valued clocks that
all advance at rate one. A finite region partition supports an exact
time-abstract reachability construction. Zones instead retain convex sets
of valuations and admit efficient difference-bound operations; their raw
reachable graph need not be finite. This treatment develops those two
constructions separately, following the clock and cell representation of
\textcite[\S\S3.1--3.4,4.1--4.3,5.5]{AlurDill1994TimedOriginal}
and the invariant/zone presentation of
\textcite[\S\S29.2--29.3,29.7.1--29.7.3]{BouyerEtAl2018TimedChapter}.
The result here is finite location reachability, not a replacement for their
infinite-word acceptance, temporal-logic, weighted, or game treatments.

\section{Clocks and the two kinds of transition}
\label{sec:p09-timed-clock-semantics}

\begin{hbdefinition}{def:p09-timed-model}{The specified timed-automaton class}
A clock valuation on a finite set \(C\) is a map
\(v:C\to\mathbb R_{\geq0}\). Write \(0_C\) for the all-zero valuation,
\((v+d)(x)=v(x)+d\), and \(v[r]\) for simultaneously resetting every
clock in \(r\subseteq C\) to zero and leaving the other clocks unchanged.
A diagonal-free clock constraint is a finite conjunction of atoms
\(x\bowtie k\), with \(k\in\mathbb Q\) and
\(\bowtie\in\{<,\leq,=,\geq,>\}\). The empty conjunction is true.
An extended difference constraint additionally permits \(x-y\bowtie k\).

A timed automaton is \(A=(L,\ell_0,C,\Sigma,I,E)\), with finite locations,
finite action alphabet, finite edge set, invariant \(I(\ell)\) at each
location, and edges \(e=(\ell,g,a,r,\ell')\). Unless the diagonal extension
is explicitly invoked, guards and invariants are diagonal-free constraints.
A state is \((\ell,v)\) with \(v\models I(\ell)\). Assume
\(0_C\models I(\ell_0)\); otherwise there is no initial state in this
semantics. Delay \(d\geq0\) is permitted at \((\ell,v)\) exactly when
\(v+t\models I(\ell)\) for every \(t\in[0,d]\). An edge may then fire
when \(v+d\models g\), producing \((\ell',(v+d)[r])\), which must satisfy
\(I(\ell')\). Write this combined delay/action step as
\((\ell,v)\xrightarrow{d,a}(\ell',(v+d)[r])\).

The original Alur--Dill transition tables have no location invariants and
use Boolean combinations of unary rational comparisons. The invariant
class above is the separately declared later convention. A finite Boolean
guard can be split into finitely many conjunctive guarded edges; arbitrary
Boolean invariants are not silently added, because the endpoint argument
below uses convexity. The chosen finite-reachability semantics allows zero
delays; this does not identify it with a convention requiring strictly
increasing timestamps of an infinite timed word.
\end{hbdefinition}

\begin{hblemma}{thm:p09-timed-rescaling}{proved=annex:p09-timed-rescaling}{Rational constants may be scaled to integers}
Let \(q>0\) be an integer divisible by all denominators of the finitely
many rational constants in \(A\). Multiply every constant, clock value,
and delay by \(q\). This gives a bijection between finite runs of the
original and scaled automata, preserving control locations, actions,
resets, invariant satisfaction, and time divergence of infinite runs when
such runs are otherwise defined. The following region construction is
therefore stated for integer constants. This rescaling changes the unit of
time, not the event order or the rates relative to that new unit.
\end{hblemma}

\begin{hbexample}{def:p09-timed-two-clock-trace}{Independent resets, one global rate}
Alur--Dill's two-clock example constrains an \(a,c\) delay to be less than
one and a \(b,d\) delay to be greater than two. Reset \(x\) on \(a\),
reset \(y\) on \(b\), test \(x<1\) on \(c\), and test \(y>2\) on
\(d\). For timestamps \(2,2.7,2.8,5\), the post-event valuations are
\((0,2),(0.7,0),(0.8,0.1),(3,2.3)\). Both tests pass. Different reset
times measure different elapsed intervals, but both clocks still share the
same rate-one time; these are not drifting local distributed clocks
\cite[exs.~3.5,3.9, pp.~192--195]{AlurDill1994TimedOriginal}.
\end{hbexample}

\section{The bounded-cell region partition}
\label{sec:p09-timed-region-partition}

\begin{hbdefinition}{def:p09-timed-region-equivalence}{Region descriptors with an explicit over-limit bit}
For every clock choose an integer \(M_x\geq0\) at least as large as the
absolute value of every unary constant tested for that clock in guards or
invariants; an otherwise unused clock may use \(M_x=0\). Put
\(n=|C|\), and write \(\langle t\rangle=t-\lfloor t\rfloor\).
Valuations \(v,w\) satisfy \(v\sim_M w\) iff all these requirements hold:

\begin{enumerate}
\item For each \(x\), \(v(x)>M_x\) iff \(w(x)>M_x\).
\item If \(v(x)\leq M_x\), their integer parts agree and
      \(\langle v(x)\rangle=0\) iff \(\langle w(x)\rangle=0\).
\item For bounded clocks \(x,y\),
      \(\langle v(x)\rangle\leq\langle v(y)\rangle\) iff
      \(\langle w(x)\rangle\leq\langle w(y)\rangle\).
\end{enumerate}

Thus each clock occupies exactly one member of the cell family
\[
 \bigl\{\{k\}:k\in\mathbb Z,\ 0\leq k\leq M_x\bigr\}
 \ \cup\ \bigl\{(k,k+1):k\in\mathbb Z,\ 0\leq k<M_x\bigr\}
 \ \cup\ \bigl\{(M_x,\infty)\bigr\},
\]
which for \(M_x=0\) consists exactly of \(\{0\}\) and
\((0,\infty)\),
and the descriptor records the weak ordering, including ties, of the
bounded nonzero fractional parts. A \emph{region} is an equivalence class
of this explicitly typed partition. This is the point/open-cell
descriptor presentation used for the construction
\cite[\S4.2]{AlurDill1994TimedOriginal}; the maximum integer point and the
over-limit cell are distinct. Extend equivalence to states only at the same
location. Guard and invariant constants must be included in the bounds.
\end{hbdefinition}

\begin{hbtheorem}{thm:p09-timed-region-properties}{proved=annex:p09-timed-region-properties}{Finite regions, stable guards/resets, and matched time elapse}
The number \(R_M\) of regions satisfies
\[
 R_M\leq n!\,2^n\prod_{x\in C}(2M_x+2).
\]
This includes \(n=0\), with empty product and \(0!\) equal to one.
Equivalent valuations agree on every allowed unary constraint, and remain
equivalent after any simultaneous reset. If \(v\sim_M w\), there is a
continuous strictly increasing bijection
\(h:[0,\infty)\to[0,\infty)\), with \(h(0)=0\), such that
\(v+t\sim_M w+h(t)\) for every \(t\geq0\).
Consequently state-region equivalence is a time-abstract bisimulation of
the specified combined-step semantics: it matches the same action/edge but
may change the numerical delay. It is not an equality-of-delays bisimulation.
\end{hbtheorem}

\begin{hbdefinition}{def:p09-timed-region-graph}{The exact quotient graph}
Vertices are the state regions satisfying their location invariant, with
initial vertex \((\ell_0,[0_C])\). An edge labelled \(a\) joins two
vertices iff some representatives make an allowed combined \(d,a\) step.
The target set is all vertices whose location lies in a fixed
\(F\subseteq L\). The witnesses in this definition are existential,
not a choice of independently incompatible representatives at every step.
The bisimulation theorem is what justifies splicing quotient edges.
\end{hbdefinition}

\begin{hbtheorem}{thm:p09-timed-region-reachability}{proved=annex:p09-timed-region-reachability}{Finite exact time-abstract reachability}
A location in \(F\) is reachable by a finite run iff a target vertex is
reachable in the region graph. There are at most \(|L|R_M\) vertices.
The graph can be constructed effectively from its cell descriptors by
enumerating their finite ordered integer-hit events, testing guards and
invariants, and applying resets. Thus finite location reachability is
decidable. This statement does not impose a non-Zeno acceptance condition
on infinite paths.
\end{hbtheorem}

\begin{hbexample}{def:p09-timed-28-regions}{The original two-clock region example}
For \(M_x=2,M_y=1\), the Alur--Dill example has 28 regions
\cite[ex.~4.4, fig.~9]{AlurDill1994TimedOriginal}. In the bounded rectangle
there are six integer points, seven horizontal/vertical open edges, two
fractional-equality diagonals, and four open two-dimensional cells: 19
regions. Outside that rectangle there are five one-dimensional rays and
four open cells, giving nine more. Hence the source's dimensional counts
are six point regions, fourteen one-dimensional regions, and eight
two-dimensional regions. For example \((0.3,0.7)\) is in \(0<x<y<1\).
The valuations \((1.3,0.7)\) and \((1.7,0.3)\) have the same integer
parts but lie on opposite sides of \(y=x-1\): the first has
\(\langle x\rangle<\langle y\rangle\), the second the reverse.
Integer parts alone would miss this difference.
\end{hbexample}

\begin{hbdefinition}{def:p09-timed-diagonal-refinement}{A finite diagonal-sensitive refinement}
For an automaton using extended difference constraints, let \(\mathcal D\)
be its finite set of tested diagonal atoms. Choose a common integer
\(M\geq0\) bounding the absolute values of all unary and diagonal
constants, and use \(M_x=M\) in the ordinary partition. Define
\(v\sim_{M,\mathcal D}w\) iff \(v\sim_Mw\) and \(v,w\) agree on
the truth value of every atom in \(\mathcal D\). These extra bits are
part of the descriptor, not assertions that differences are bounded by
the individual clock bounds.
\end{hbdefinition}

\begin{hbtheorem}{thm:p09-timed-diagonal-refinement}{proved=annex:p09-timed-diagonal-refinement}{The diagonal refinement preserves finite reachability}
The refined partition has at most \(2^{|\mathcal D|}R_M\) regions. Guard,
invariant, reset and matched-delay properties, and the quotient finite-run
reachability equivalence, hold for this refinement. Its feasible descriptors
and edges are effectively obtainable using the difference-constraint
feasibility construction below. The \emph{unrefined} ordinary partition is
not sufficient for arbitrary diagonal guards.
\end{hbtheorem}

\begin{hbexample}{def:p09-timed-diagonal-failure}{Why the extra diagonal bits matter}
Start with \(x=y=0\), wait until \(y=1\), reset \(x\), and enter \(p\).
In \(p\), allow arbitrary delay but no further reset. Every reachable
valuation satisfies \(y=x+1\). An outgoing guard \(x-y>0\) is never
enabled. With ordinary bounds \(M_x=M_y=1\), reachable \((2,3)\) and
unreachable \((3,2)\) are in the same over-limit region, but disagree on
that guard. Ordinary region saturation could therefore create a false
target path. The refined truth bit separates them; exact zones retain the
difference \(y-x=1\) as well.
\end{hbexample}

\section{Zones and strict difference-bound matrices}
\label{sec:p09-timed-dbm}

\begin{hbdefinition}{def:p09-timed-dbm}{The DBM constraint class and conventions}
Index the clocks by \(1,\ldots,n\) and add a fixed coordinate \(x_0=0\).
A bound is \((c,\triangleleft)\), where \(c\in\mathbb Q\) and
\(\triangleleft\in\{<,\leq\}\), or \(\infty\). Bound order
\(\preceq_b\) means no weaker: smaller constants are tighter, and
\((c,<)\prec_b(c,\leq)\). Bound addition \(\oplus_b\) adds constants
and makes the result strict iff either summand is strict; infinity is
absorbing. A DBM \(D\) represents
\[
 Z(D)=\{v\in\mathbb R_{\geq0}^n:
               v_i-v_j\mathrel{\triangleleft_{ij}}c_{ij}
               \text{ for every finite }D_{ij}\},\qquad v_0=0.
\]
Always insert \(D_{ii}\preceq_b(0,\leq)\) and the nonnegativity bounds
\(D_{0i}\preceq_b(0,\leq)\). Equalities use two nonstrict bounds;
lower bounds reverse the two indices and negate the constant. Empty
zones have a separate empty marker propagated by every operation.

The weighted graph convention used in the feasibility proof is an edge
\(j\to i\) of bound \(D_{ij}\); it implements the displayed
\(v_i-v_j\) constraint. This is the transpose orientation of the
constraint-labelled arrows in the chapter's illustration, not a change
of constraint meaning. Shortest-path closure is computed by the
Floyd--Warshall updates
\(D_{ij}\leftarrow\min_b(D_{ij},D_{ik}\oplus_bD_{kj})\), for each
\(k=0,\ldots,n\). A closed nonempty matrix is called canonical here.
The class is rational difference constraints only: not arbitrary linear
polyhedra, clock sums, or parametric bounds.
\end{hbdefinition}

\begin{hbtheorem}{thm:p09-timed-dbm-closure}{proved=annex:p09-timed-dbm-closure}{Closure and strict-bound consistency}
Closure preserves \(Z(D)\). The zone is empty iff closure has some
\(D_{ii}\prec_b(0,\leq)\): a negative constant or a zero strict bound.
Otherwise it has a rational valuation. Closure uses \(O((n+1)^3)\)
bound operations; that count does not treat unbounded rational bit
arithmetic as unit-cost in a complexity claim.
\end{hbtheorem}

\begin{hbtheorem}{thm:p09-timed-dbm-inclusion}{proved=annex:p09-timed-dbm-inclusion}{Exact inclusion of canonical zones}
For canonical nonempty DBMs \(D,E\),
\(Z(D)\subseteq Z(E)\) iff \(D_{ij}\preceq_bE_{ij}\) for all indices.
The empty-marker cases are decided first. Comparing constants without
their strictness flags is not an inclusion test.
\end{hbtheorem}

\begin{hbtheorem}{thm:p09-timed-dbm-operations}{proved=annex:p09-timed-dbm-operations}{Exact intersection, reset, and time elapse}
For a nonempty canonical matrix \(D\):
\begin{enumerate}
\item Intersection with another zone or conjunctive guard is represented
      by entrywise minimum followed by closure and an emptiness check.
\item To reset \(r\), put \(\rho(i)=0\) for reset indices and
      \(\rho(i)=i\) otherwise, including \(\rho(0)=0\). The matrix
      \(D^r_{ij}=D_{\rho(i),\rho(j)}\) represents exactly \(Z(D)[r]\).
\item To allow arbitrary common delay, replace \(D_{i0}\) by infinity
      for every \(i>0\), retaining all other entries. The result represents
      exactly \(Z(D)^\uparrow=\{v+d:v\in Z(D),d\geq0\}\).
\end{enumerate}
Zero resets leave the set unchanged. Empty inputs remain empty; resetting
coordinates in an inconsistent raw matrix must not erase its inconsistency.
Canonical closure is an actual precondition for the simple reset/time-up
formulas, not an optional formatting step.
\end{hbtheorem}

\begin{hbexample}{def:p09-timed-dbm-calculation}{One exact DBM calculation}
Take \(Z=\{0\leq x\leq2,\ y=x+1\}\). In order \(0,x,y\), its closed
nonstrict-bound matrix has constants
\[
 D=\begin{pmatrix}0&0&-1\\2&0&-1\\3&1&0\end{pmatrix}.
\]
Time-up removes the entries \(2,3\) in column zero and gives
\(\{x\geq0,y=x+1\}\). Intersecting the original zone with \(x<1\)
and closing also derives \(y<2\). Resetting \(x\) then yields
\(\{x=0,1\leq y<2\}\), with the strict upper bound retained.
Resetting \(x\) in the original \(Z\) instead gives
\(\{x=0,1\leq y\leq3\}\): reset forgets its old value; it does not
substitute zero into every old difference equation.
\end{hbexample}

\begin{hbexample}{def:p09-timed-clock-domain-check}{Nonnegativity is part of the zone}
Among the constraints in the chapter's Figure~12 are \(x_1\leq3\) and
\(x_1-x_2\geq4\) \cite[p.~1025]{BouyerEtAl2018TimedChapter}.
As constraints on nonnegative clocks, these already imply
\(4\leq x_1\leq3\), so their clock zone is empty. A weighted graph
without the implicit nonnegativity rows can still describe a nonempty
real-coordinate polyhedron. Our DBMs always include those rows and do
not mistake that polyhedron for a nonempty clock zone.
\end{hbexample}

\section{Exact posts and a genuinely finite normalization}
\label{sec:p09-timed-zone-reachability}

\begin{hbdefinition}{def:p09-timed-zone-post}{Symbolic combined-step post}
For \(Z\subseteq\llbracket I(\ell)\rrbracket\), define
\[
 \operatorname{Post}_e(Z)=
    \bigl(Z^\uparrow\cap\llbracket I(\ell)\rrbracket
                         \cap\llbracket g\rrbracket\bigr)[r]
                         \cap\llbracket I(\ell')\rrbracket.
\]
A symbolic path starts at the singleton zero zone, and uses this post for
each selected edge, discarding empty results. Keep separate zones or their
exact union; replacing that union by its convex hull is not this algorithm.
\end{hbdefinition}

\begin{hbtheorem}{thm:p09-timed-zone-post}{proved=annex:p09-timed-zone-post}{Zone posts describe precisely the finite runs of each edge sequence}
For conjunctive difference-constraint invariants and guards, the symbolic
post is exactly the set of combined-step successors. Consequently the
zone associated with a finite edge sequence is exactly the set of valuations
reached by its concrete runs. Every needed post is computed by the proved
DBM operations. This gives exact symbolic finite-path reachability but
does not, without a separate finiteness argument, give a terminating search
over every zone path.
\end{hbtheorem}

\begin{hbexample}{def:p09-timed-infinite-zones}{An infinite raw zone graph}
Use clocks \(x,y\), one location with invariant \(x\leq1\), and a loop
guarded by \(x=1\) which resets only \(x\). After its \(k\)-th traversal
the exact zone is \(\{x=0,y=k\}\), for every \(k\geq0\). These are
infinitely many distinct DBMs, although finite regions abstract them:
\(M_x=1,M_y=0\) suffices, and all \(k>0\) share the same over-limit
\(y\) descriptor. Each iteration takes one time unit, so this failure
of zone-graph finiteness is not caused by Zeno behavior.
\end{hbexample}

\begin{hbexample}{def:p09-timed-hull-failure}{A convex hull can create an error path}
From the zero state, two edges into \(p\) use guards \(x=0\) and \(x=2\),
respectively, and reset an auxiliary clock \(z\). At \(p\), invariant
\(z\leq0\) forbids positive delay. The exact reachable set there is
\(\{(x,z)=(0,0),(2,0)\}\). An outgoing guard \(x=1\) is impossible.
Its convex hull contains \((1,0)\), falsely enabling that edge. A union
of zones is therefore not in general one zone or an exact DBM hull.
\end{hbexample}

\begin{hbdefinition}{def:p09-timed-region-saturation}{Region-union normalization}
For the applicable ordinary or diagonal-refined partition, put
\[
 \alpha(Z)=\bigcup\{R:R\cap Z\ne\varnothing\}.
\]
This is a finite union of regions, not generally a single zone. Normalize
symbolic successors using \(\alpha\), representing unions as explicit
region sets. This mathematical normalization is distinct from practical
maximum-bound/LU DBM extrapolations; no unproved extrapolation is substituted
for it \cite[\S29.7.1]{BouyerEtAl2018TimedChapter}.
\end{hbdefinition}

\begin{hbtheorem}{thm:p09-timed-finite-normalization}{proved=annex:p09-timed-finite-normalization}{Finite normalized search with exact location reachability}
For the appropriate bisimulation partition,
\[
 \alpha\bigl(\operatorname{Post}_e(\alpha(Z))\bigr)
       =\alpha\bigl(\operatorname{Post}_e(Z)\bigr).
\]
Normalized symbolic search preserves finite location reachability. There
are at most \(|L|2^{R}\) normalized symbolic states, where \(R\) is the
applicable region bound. Alternatively accumulate a set of reachable
regions at each location; at most \(|L|R\) new region entries can be added.
Guard satisfaction and every finite symbolic path lift to concrete runs by
bisimulation, not by treating arbitrary abstract valuations as reachable.
Ordinary normalization in the presence of unchecked diagonal guards does
not satisfy these hypotheses.
\end{hbtheorem}

\begin{hbwarning}{def:p09-timed-acceptance-boundary}{Finite reachability does not establish infinite acceptance}
A loop with guard and invariant \(x<1\), never resetting \(x\), can have
infinitely many traversals at increasing times converging to one, but no
time-divergent such run. A control/region cycle alone is not an accepting
progressing timed run. Alur--Dill's original infinite-word construction
adds progressive-run conditions; the 2018 chapter separately defines
time-divergent runs and non-Zeno extendibility. Those conditions are not
conclusions of the finite reachability theorem. Likewise, time-abstract
matching does not preserve exact timestamps, drifted clocks, arbitrary
clock sums, stopped clocks, general hybrid updates, parametric bounds,
language inclusion, or controller winning conditions. Each requires its
own model and proof. The CTMC residence-time semantics of
\cref{sec:p09-csl-endpoints} is probabilistic timing, not this clock-guard
transition semantics.
\end{hbwarning}
